% GENERATED FILE. Edit canonical lessons, then run npm run generate:books.
% canonical-source-hash: fnv1a64:e0f13b2b7ab7082d
% canonical-lessons: ZH-C350-shishi, ZH-C350-zhengming, ZH-C350-genju, ZH-C350-panduan, ZH-C350-queshi

\chapter{Facts and Grounds}
\label{ch:zh-facts-and-grounds}
\hlchaptermodality{350}

\begin{chapteropening}
\textbf{By the end of this chapter:} \emph{I can say what the facts are, what a claim rests on and that it is hard to judge.}

Complete the last lesson of chapter 350: I can say what the facts are, what a claim rests on and that it is hard to judge.
\end{chapteropening}

\section[shìshí]{shìshí — fact}

\label{lesson:ZH-C350-shishi}



\begin{quote}
Before the new one: say the Mandarin for the way something is done, then the Mandarin for to argue in favour of.
\end{quote}

\subsection*{You'll want to know: shìshí}
\emph{shìshí} — \textquotedblleft{}fact\textquotedblright{}.

\begin{quote}
\emph{wǒ juéde tā shuō de shì shìshí} — I think what he says is true.
\end{quote}

Its characters come later; for now it is a word for the ear.

\begin{cousinweb}
\emph{shìshí} is built from \zh{事} 'a matter' and \emph{shí}, both inherited from Middle Chinese, \emph{dzriH} and \emph{zyit} (Guangyun, in Baxter's transcription), which glosses the second 'full' and 'sincere'.
\end{cousinweb}

\begin{grammarlens}[title={What you've built}]
One more word for giving reasons.
\end{grammarlens}

\subsection*{Your turn}
\begin{itemize}
\raggedright
  \item \emph{Say it:} \emph{shìshí}
  \item \emph{Say it:} \emph{shìshí}, once more
  \item \emph{Say it:} \emph{wǒ juéde tā shuō de shì shìshí} — I think what he says is true.
  \item \emph{Recall:} say the Mandarin for the way something is done, then the Mandarin for to argue in favour of, then say \emph{shìshí} again
\end{itemize}

\begin{tcolorbox}[breakable,colback=teal!4,colframe=teal!35!black,title={Before you move on}]
What does shìshí mean? (\textquotedblleft{}Fact\textquotedblright{}.) Say it once more.
\end{tcolorbox}
\section[zhèngmíng]{zhèngmíng — to prove}

\label{lesson:ZH-C350-zhengming}



\begin{quote}
Before the new one: say the Mandarin for to argue in favour of, then the Mandarin for a fact.
\end{quote}

\subsection*{You'll want to know: zhèngmíng}
\emph{zhèngmíng} — \textquotedblleft{}to prove\textquotedblright{}.

\begin{quote}
\emph{nǐ néng zhèngmíng nǐ shuō de huà ma} — Can you prove what you say?
\end{quote}

Its characters come later; for now it is a word for the ear.

\begin{cousinweb}
\emph{zhèngmíng} is built from \emph{zhèng} and \zh{明} 'bright', both inherited from Middle Chinese, \emph{tsyingH} and \emph{mjaeng} (Guangyun, in Baxter's transcription), which glosses the first 'to verify' and the second 'bright'.
\end{cousinweb}

\begin{grammarlens}[title={What you've built}]
One more word for giving reasons.
\end{grammarlens}

\subsection*{Your turn}
\begin{itemize}
\raggedright
  \item \emph{Say it:} \emph{zhèngmíng}
  \item \emph{Say it:} \emph{zhèngmíng}, once more
  \item \emph{Say it:} \emph{nǐ néng zhèngmíng nǐ shuō de huà ma} — Can you prove what you say?
  \item \emph{Recall:} say the Mandarin for to argue in favour of, then the Mandarin for a fact, then say \emph{zhèngmíng} again
\end{itemize}

\begin{tcolorbox}[breakable,colback=teal!4,colframe=teal!35!black,title={Before you move on}]
What does zhèngmíng mean? (\textquotedblleft{}To prove\textquotedblright{}.) Say it once more.
\end{tcolorbox}
\section[gēnjù]{gēnjù — basis; what something rests on}

\label{lesson:ZH-C350-genju}



\begin{quote}
Before the new one: say the Mandarin for a fact, then the Mandarin for to prove.
\end{quote}

\subsection*{You'll want to know: gēnjù}
\emph{gēnjù} — \textquotedblleft{}basis; what something rests on\textquotedblright{}.

\begin{quote}
\emph{yīnwèi méiyǒu gēnjù suǒyǐ wǒ bù xiāngxìn} — There's nothing to back it up, so I don't believe it.
\end{quote}

Its characters come later; for now it is a word for the ear.

\begin{cousinweb}
\emph{gēnjù} is built from \emph{gēn} and \emph{jù}, both inherited from Middle Chinese, \emph{kon} and \emph{kjoH} (Guangyun, in Baxter's transcription), which glosses the first 'a root' and the second 'to rely on'.
\end{cousinweb}

\begin{grammarlens}[title={What you've built}]
One more word for giving reasons.
\end{grammarlens}

\subsection*{Your turn}
\begin{itemize}
\raggedright
  \item \emph{Say it:} \emph{gēnjù}
  \item \emph{Say it:} \emph{gēnjù}, once more
  \item \emph{Say it:} \emph{yīnwèi méiyǒu gēnjù suǒyǐ wǒ bù xiāngxìn} — There's nothing to back it up, so I don't believe it.
  \item \emph{Recall:} say the Mandarin for a fact, then the Mandarin for to prove, then say \emph{gēnjù} again
\end{itemize}

\begin{tcolorbox}[breakable,colback=teal!4,colframe=teal!35!black,title={Before you move on}]
What does gēnjù mean? (\textquotedblleft{}Basis; what something rests on\textquotedblright{}.) Say it once more.
\end{tcolorbox}
\section[pànduàn]{pànduàn — to make a judgement; to judge}

\label{lesson:ZH-C350-panduan}



\begin{quote}
Before the new one: say the Mandarin for to prove, then the Mandarin for a basis for a claim.
\end{quote}

\subsection*{You'll want to know: pànduàn}
\emph{pànduàn} — \textquotedblleft{}to make a judgement; to judge\textquotedblright{}.

\begin{quote}
\emph{wǒ bù néng pànduàn tā shuō de shì bu shì shìshí} — I can't judge whether what he says is true.
\end{quote}

Its characters come later; for now it is a word for the ear.

\begin{cousinweb}
\emph{pànduàn} is built from \emph{pàn} and \emph{duàn}, both inherited from Middle Chinese, \emph{phanH} and \emph{dwanX} (Guangyun, in Baxter's transcription), which glosses the first 'to split' and 'to divide', and the second 'to cut off'; it also lists the second in the departing tone, glossed 'to decide'.
\end{cousinweb}

\begin{grammarlens}[title={What you've built}]
One more word for giving reasons.
\end{grammarlens}

\subsection*{Your turn}
\begin{itemize}
\raggedright
  \item \emph{Say it:} \emph{pànduàn}
  \item \emph{Say it:} \emph{pànduàn}, once more
  \item \emph{Say it:} \emph{wǒ bù néng pànduàn tā shuō de shì bu shì shìshí} — I can't judge whether what he says is true.
  \item \emph{Recall:} say the Mandarin for to prove, then the Mandarin for a basis for a claim, then say \emph{pànduàn} again
\end{itemize}

\begin{tcolorbox}[breakable,colback=teal!4,colframe=teal!35!black,title={Before you move on}]
What does pànduàn mean? (\textquotedblleft{}To make a judgement; to judge\textquotedblright{}.) Say it once more.
\end{tcolorbox}
\section[quèshí]{quèshí — indeed; it is true that}

\label{lesson:ZH-C350-queshi}



\begin{quote}
Before the new one: say the Mandarin for a basis for a claim, then the Mandarin for to make a judgement.
\end{quote}

\subsection*{You'll want to know: quèshí}
\emph{quèshí} — \textquotedblleft{}indeed; it is true that\textquotedblright{}.

\begin{quote}
\emph{tā shuō de quèshí shì shìshí} — What he said really is a fact.
\end{quote}

Its characters come later; for now it is a word for the ear.

\begin{cousinweb}
\emph{quèshí} is built from \emph{què} and \emph{shí}, both inherited from Middle Chinese, \emph{khaewk} and \emph{zyit} (Guangyun, in Baxter's transcription), which glosses the first 'firm' and the second 'full' and 'sincere'.
\end{cousinweb}

\begin{grammarlens}[title={What you've built}]
One more word for giving reasons.
\end{grammarlens}

\subsection*{Your turn}
\begin{itemize}
\raggedright
  \item \emph{Say it:} \emph{quèshí}
  \item \emph{Say it:} \emph{quèshí}, once more
  \item \emph{Say it:} \emph{tā shuō de quèshí shì shìshí} — What he said really is a fact.
  \item \emph{Recall:} say the Mandarin for a basis for a claim, then the Mandarin for to make a judgement, then say \emph{quèshí} again
\end{itemize}

\begin{tcolorbox}[breakable,colback=teal!4,colframe=teal!35!black,title={Before you move on}]
What does quèshí mean? (\textquotedblleft{}Indeed; it is true that\textquotedblright{}.) Say it once more.
\end{tcolorbox}
